% Pursuing Stacks (cotes 134-2 to 134-8): archive cote 134-7, site page 27 (PDF page 28), typescript page 469.
% typescript. Draft compared: ps6.tex:1253-1313.
% AI-transcribed by Claude Opus 5.5 (claude-opus-5-5, adaptive thinking, effort high), Message Batches API, batch ps3 (msgbatch_01XdFmuVVBAcgc7EFxdtCqsj), 2026-11-12; not yet checked by eye.
\scanpage{134-7/27}
\mspage{469}

because of the possibility of operation of \hand{$\pi$}$_1$ upon the \hand{$\pi$}$_i$'s. Maybe,
when trying to modelize usual homo\xout{b}to\xout{x}py types by complexes of unipotent
bundles over $\underline{Z}$, we should restrict to homotopy types which are not only
\xout{zer}0-connected and have abelian \hand{$\pi$}$_1$ , but moreover with \hand{$\pi$}$_1$ operating
trivially (or for the very least, in a unipotent way) upon the higher
\hand{$\pi$}$_i$'s. At any rate, as soon as \hand{$\pi$}$_1$ operates non trivially (on itself, or
on the higher \hand{$\pi$}$_i$'s ) there will presumably be two non equivalent ways for
defining soft Postnikov dévissage, corresponding to the two standard
descending commutator group series in a discrete group G. The more relevant
\xout{xxx} in view of unipotent schematization would seem to be the ``nilpotent''
one.

Restricting for simplicity to the 1-connected case \unsure{$X_\bullet$|$X_0$}$=X_1=e$ , I would
expect soft Postnikov dévissage to be the key \unsure{\hand{f}or} an understanding as well
of the behaviour of the \hand{$\pi$}$_i(X_*)$ modules with respect to ring extension,
as \xout{xxx}\inserted[typed]{of} the full relationship between these invarian ts, and the homology
invarian ts $H_i(X_*)$.

$*$ \quad $*$

$*$

I still should have a look upon complexes $X_*$ in $U(k)^\bullet$\hand{/}satisfying
\hand{(}as always in this game\hand{)} $X_0=e$ , but not necessarily $X_1 = e$. Even when
we make the Kan assumption (\deleted{including}\inserted{plus} ``\uline{smoothness}'' of the components,
by which I mean that their k-linearization is projective), I do'nt feel
too sure yet if they fit into a good formalism, for instance \hand{(}when $k=\underline{Z}$\hand{)}
if they satisfy the ``linearization theorem'' (quasiisomorphy for the two
maps in (w) p.454). If we start for instance with a group object\inserted{G} of U(k)
and let $X_*$ be the corresponding semisimplicial complex in $U(k)^\bullet$ , then
we get an isomorphism
\[ \text{(E)} \qquad \hand{$\pi$}_1(X_*) \overset{\text{def}}{=} \hand{$\pi$}_1(X_*(k)) \simeq G(k) \quad , \]
which shows us that \hand{$\pi$}$_1(X_*)$ need not be abelian even when $k=\underline{Z}$ , \deleted{all we \xout{sa}}
\deleted{can safely say,} \over{i}{I}f we assume that $X_1 = G$ is ``of finite presentation''
(namely the \inserted[typed]{projective} k-module which describes $X_1$ is of finite presentation), \deleted{is}\inserted{or}
what amounts to the same, representable by an actual (group) \unsure{scheme|shheme}, \inserted{it}
is \inserted{true however} that G and hence \hand{$\pi$}$_1 = G(k)$ is \uline{nilpotent} (this holds for any k ).
It looks an intriguing question whether \hand{$\pi$}$_1$ is nilpotent under the only
assumption that $X_*$ is a \inserted{smooth} Kan complex with $X_1$ a scheme (wi\over{th}{th}out assuming
any more X comes from a group object). At any rate, it follows from

\margin{27}
